\section{Spektrum suatu gelanggang}

Sekarang kita akan mengaitkan setiap gelanggang komutatif dengan suatu ruang
topologis $\spec R$
secara funktorial.
Artinya, akan terdapat suatu funktor kontravarian
\[\spec:  \mathbf{CRing} \to \mathbf{Top}  \]
dengan $\mathbf{Top}$ kategori ruang topologis.
Konstruksi ini menjadi dasar geometri aljabar
teoretis-skema dan akan sering digunakan selanjutnya. 

Pengamatan yang memotivasinya adalah sebagai berikut. Jika $k$ suatu
medan tertutup secara aljabar, maka ideal-ideal maksimal dalam $k[x_1, \dots, x_n]$ berbentuk
$(x_1-a_1, \dots, x_n  - a_n)$ untuk $(a_1, \dots, a_n) \in k[x_1, \dots, x_n]$.
Inilah Nullstellensatz, yang belum kita buktikan. Dengan demikian, kita dapat
mengidentifikasi ideal-ideal maksimal dalam gelanggang polinomial dengan ruang $k^n$.
Jika $I \subset k[x_1, \dots, x_n]$ suatu ideal, maka ideal-ideal maksimal dalam
$k[x_1,\dots,x_n]$ berkorespondensi dengan titik-titik tempat semua elemen $I$ lenyap. Lihat 
\rref{twovarpoly} untuk penjelasan yang lebih terperinci. Jadi, geometri aljabar afin klasik
mengkaji himpunan ideal maksimal dalam suatu aljabar yang terbangkitkan
hingga di atas medan tertutup secara aljabar.

$\mathrm{Spec}$ suatu gelanggang merupakan generalisasi konstruksi ini.
Secara umum, lebih wajar untuk 
menggunakan semua ideal prima, bukan hanya ideal maksimal.
\subsection{Definisi dan contoh}

Kita mulai dengan mendefinisikan $\spec$ sebagai suatu himpunan. Selanjutnya kita akan
mengkonstruksi
topologi Zariski, lalu sifat funktorialnya.
\begin{definition} 
Misalkan $R$ suatu gelanggang komutatif. \textbf{Spektrum} $R$,
yang dinotasikan dengan $\spec R$, adalah
himpunan ideal prima dari $R$.
\end{definition} 

Sekarang kita akan menjadikan $\spec R$ suatu ruang topologis. Pertama-tama, kita
deskripsikan suatu
koleksi himpunan yang nantinya menjadi himpunan tertutup. 
Jika $I \subset R$ suatu ideal, tetapkan
\[ V(I) = \left\{\mathfrak{p}: \mathfrak{p} \supset I\right\}
\subset \spec R.
\]

\begin{proposition} 
Terdapat suatu topologi pada $\spec R$ sedemikian sehingga subhimpunan-subhimpunan tertutupnya
berbentuk
$V(I)$ untuk suatu ideal $I \subset R$.
\end{proposition} 

\begin{proof} 
Kita memang harus memeriksa aksioma-aksioma yang lazim bagi suatu topologi:
\begin{enumerate}
\item $\emptyset = V((1))$ karena tidak ada ideal prima yang memuat $1$. Jadi
$\emptyset$ tertutup.
\item $\spec R = V((0))$ karena setiap ideal memuat nol. Jadi
$\spec R$
tertutup.
\item Kita tunjukkan bahwa himpunan-himpunan tertutup stabil terhadap irisan.
Misalkan
$K_{\alpha} = V(I_{\alpha})$ subhimpunan tertutup dari $\spec R$ dengan
$\alpha$
merentang suatu himpunan indeks. Misalkan $I
= \sum I_{\alpha}$. Maka 
\[ V(I) = \bigcap K_{\alpha} = \bigcap V(I_{\alpha}),  \]
yang berlaku karena $I$ adalah ideal terkecil yang memuat setiap
$I_{\alpha}$;
jadi suatu ideal prima memuat setiap $I_{\alpha}$ jika dan hanya jika ia memuat $I$.  
\item Gabungan dua himpunan tertutup bersifat tertutup. Memang, jika $K,K'
\subset \spec
R$ tertutup, kita tunjukkan bahwa $K \cup K'$ tertutup. Tuliskan $K= V(I), K' =
V(I')$. Maka
kita menyatakan bahwa 
\[ K \cup K'  = V(II').  \]
Di sini, seperti biasa, $II'$ adalah ideal yang dibangkitkan oleh hasil kali $ii', i \in I, i'
\in I'$. Jika
$\mathfrak{p}$ \textbf{prima} dan memuat $II'$, maka ia harus
memuat salah satu dari $I$, $I'$;
hal ini menyiratkan persamaan yang ditampilkan di atas dan dengan demikian
hasilnya.
\end{enumerate}
\end{proof} 
\begin{definition} 
Topologi pada $\spec R$ yang didefinisikan di atas disebut
\textbf{topologi
Zariski}. Dengan topologi ini, $\spec R$ kini merupakan suatu ruang topologis.
\end{definition} 

\begin{exercise} 
Apakah $\spec$ dari gelanggang nol?
\end{exercise} 

Untuk melihat geometri konstruksi ini, mari kita kerjakan
beberapa contoh.

\begin{example} 
Misalkan $R = \mathbb{Z}$, dan tinjau $\spec \mathbb{Z}$. Maka
setiap ideal prima dibangkitkan oleh satu elemen, karena
$\mathbb{Z}$ adalah suatu PID. Kita memiliki $\spec \mathbb{Z} = \{(0)\}
\cup \bigcup_{p \
\mathrm{prima}} \{ (p)\}$. Gambarannya adalah bahwa kita memiliki semua bilangan
prima yang lazim, $(2), (3),
(5), \dots, $ lalu suatu titik khusus $(0)$.

Sekarang mari kita deskripsikan subhimpunan-subhimpunan tertutupnya. Himpunan-himpunan ini berbentuk
$V(I)$ dengan $I
\subset \mathbb{Z}$ suatu ideal, sehingga $I = (n)$ untuk suatu $n \in
\mathbb{Z}$.

\begin{enumerate}
\item Jika $n=0$, subhimpunan tertutupnya adalah seluruh $\spec \mathbb{Z}$.
\item Jika $n \neq 0$, maka $n$ memiliki berhingga banyak pembagi prima.
Jadi $V((n))$ terdiri
dari ideal-ideal prima yang bersesuaian dengan pembagi-pembagi prima tersebut.  
\end{enumerate}

Selain seluruh ruang, satu-satunya subhimpunan tertutup adalah subhimpunan
hingga
yang tidak memuat $(0)$.  
\end{example} 

\begin{example} \label{twovarpoly}
Misalkan $R = \mathbb{C}[x,y]$ suatu gelanggang polinomial dalam dua peubah.
Di sini kita tidak akan memberikan deskripsi lengkap tentang $\spec R$. Namun, kita akan menuliskan
beberapa
ideal prima.

\begin{enumerate}
\item Untuk setiap pasangan bilangan kompleks $s,t \in \mathbb{C}$,
koleksi polinomial
$f \in R$ sedemikian sehingga $f(s,t) = 0$ merupakan suatu ideal prima $\mathfrak{m}
_{s,t} \subset R$. Sesungguhnya,
ideal ini maksimal, karena gelanggang residunya adalah seluruh
$\mathbb{C}$. Memang,
$R/\mathfrak{m}_{s,t} \simeq \mathbb{C}$ di bawah pemetaan $f \to
f(s,t)$.

Sesungguhnya, 
\begin{theorem}
Ideal-ideal $\mathfrak{m}_{s,t}$ adalah semua ideal maksimal dalam $R$.
\end{theorem} 
Hal ini akan mengikuti dari \emph{Nullstellensatz Hilbert} yang akan
dibuktikan kemudian
(\rref{gennullstellensatz}).
\item $(0) \subset R$ merupakan ideal prima karena $R$ suatu domain. 
\item Jika $f(x,y) \in R$ suatu polinomial tak tereduksi, maka $(f)$
merupakan ideal
prima. Hal ini ekuivalen dengan faktorisasi tunggal dalam
$R$.\footnote{Suatu pembuktian mengikuti dari hasil UFD bagi gelanggang polinomial yang dikembangkan dalam Li, Bab 5; pembuktian itu tidak diulang dalam suplemen pilihan ini.}  
\end{enumerate}

Untuk menggambar $\spec R$, kita mulai dengan menggambar $\mathbb{C}^2$, yang diidentifikasi
dengan
koleksi
ideal maksimal $\mathfrak{m}_{s,t}, s, t \in \mathbb{C}$. $\spec R$ juga memiliki
titik-titik tambahan (yang tidak tertutup),
sebagaimana
dideskripsikan di atas, tetapi untuk sekarang mari kita
tinjau topologi yang diinduksi pada $\mathbb{C}^2$ sebagai subruang dari
$\spec R$.

Subhimpunan tertutup dari $\spec R$ adalah subhimpunan $V(I)$ dengan $I$
suatu ideal,
yang dibangkitkan oleh polinomial-polinomial $\left\{f_{\alpha}(x,y)\right\}$. Kita ingin
menentukan subhimpunan $\mathbb{C}^2$ yang
diinduksi oleh $V(I)$.
Dengan kata lain, kita bertanya:
\begin{quote}
Titik-titik manakah dari $\mathbb{C}^2$ (dengan $(s,t)$ diidentifikasi dengan
$\mathfrak{m}_{s,t}$) yang terletak dalam $V(I)$?
\end{quote}
Menurut definisi, kita mengetahui bahwa $(s,t)$ bersesuaian dengan suatu titik $V(I)$ jika
dan hanya jika $I
\subset \mathfrak{m}_{s,t}$.
Hal ini berlaku jika dan hanya jika semua
$f_{\alpha} $ terletak dalam $ \mathfrak{m}_{s,t}$, yaitu jika
$f_{\alpha}(s,t) =0$ untuk setiap
$\alpha$. Jadi subhimpunan tertutup dari $\mathbb{C}^2$ (dengan
topologi Zariski
yang diinduksi) \emph{tepat merupakan subhimpunan
yang dapat didefinisikan oleh persamaan polinomial}.

Topologi ini
\textbf{jauh} lebih kasar
daripada topologi biasa. Misalnya, $\left\{(z_1,z_2):
\Re(z_1) \geq 0\right\}$
tidak tertutup-Zariski. 
Topologi Zariski begitu kasar karena kita hanya memiliki data
aljabar (yakni,
polinomial, atau elemen $R$) untuk mendefinisikan topologi tersebut.
\end{example} 

\begin{exercise} 
Misalkan $R_1, R_2$ gelanggang-gelanggang komutatif. Berikan $R_1 \times R_2$ suatu
struktur gelanggang yang alami,
dan deskripsikan $\spec( R_1 \times R_2)$ dalam kaitannya dengan
$\spec R_1$ dan
$\spec R_2$.
\end{exercise} 


\begin{exercise} 
Misalkan $X$ suatu ruang Hausdorff kompak, dan $C(X)$ gelanggang fungsi
kontinu bernilai real
$X \to \mathbb{R}$. 
Ideal-ideal maksimal dalam $\spec C(X)$ berkorespondensi bijektif dengan
titik-titik $X$,
dan topologi yang diinduksi pada $X $ (sebagai subhimpunan $\spec C(X)$ dengan topologi
Zariski)
tidak lain ialah topologi biasa.
\end{exercise}

\begin{exercise}
Buktikan hasil berikut: jika $X, Y$ adalah ruang Hausdorff
kompak dan $C(X),
C(Y)$ gelanggang-gelanggang fungsi kontinu yang bersesuaian, dan jika $C(X),
C(Y)$
isomorfik sebagai aljabar-$\mathbb{R}$, maka $X$ homeomorfik dengan
$Y$.
\end{exercise} 


\subsection{Korespondensi ideal radikal--subhimpunan tertutup}

Sekarang kita kembali ke kasus gelanggang komutatif sembarang $R$.
Jika $I \subset R$, kita memperoleh suatu subhimpunan tertutup
$V(I) \subset \spec R$. Himpunan itu disebut $V(I)$ karena kita
dimaksudkan untuk
memandangnya sebagai tempat-tempat di mana elemen-elemen $I$ ``lenyap'',
sebab
elemen-elemen $R$ menyerupai ``fungsi''. Analogi ini mungkin
paling jelas
terlihat dalam contoh gelanggang polinomial di atas medan
tertutup secara aljabar,
misalnya \rref{twovarpoly} di atas.

Pemetaan dari ideal ke himpunan tertutup sangat jauh dari bersifat
injektif secara
umum, walaupun menurut definisi ia surjektif.

\begin{example} 
Jika $R = \mathbb{Z}$ dan $p$ prima, maka $I = (p), I' = (p^2)$
mendefinisikan
subhimpunan yang sama (yakni, $\left\{(p)\right\}$) dari
$\spec R$. 
\end{example} 

Sekarang kita bertanya mengapa pemetaan dari ideal ke subhimpunan
tertutup gagal
bersifat injektif. Seperti yang akan kita lihat, seluruh masalah lenyap jika
kita membatasi perhatian pada ideal-ideal
\emph{radikal}.

\begin{definition} 
Jika $I$ suatu ideal, maka \textbf{radikal} $\rad(I) $ atau $
\sqrt{I}$
didefinisikan sebagai $$\rad(I) =
\left\{x \in R: x^n \in I \ \mathrm{untuk} \ \mathrm{suatu} \ n
\right\}.$$
Suatu ideal disebut \textbf{radikal} jika sama dengan radikalnya.
(Ini
ekuivalen dengan \rref{def-radical-ideal} sebelumnya.) 
\end{definition} 

Sebelum melanjutkan, kita harus memeriksa hal berikut.
\begin{lemma} 
Jika $I$ suatu ideal, maka $\rad(I)$ juga suatu ideal.
\end{lemma} 
\begin{proof} 
Jelas bahwa $\rad(I)$ tertutup terhadap perkalian karena $I$ demikian.
Andaikan $x,y \in \rad(I)$; kita tunjukkan bahwa $x+y \in \rad(I)$. Maka $x^n,
y^n \in I$
untuk suatu $n$ (yang besar), dan dengan demikian untuk semua $n$ yang lebih besar. Ekspansi binomial kini
memberikan
\[ (x+y)^{2n} = x^{2n} + \binom{2n}{1} x^{2n-1}y + \dots +
y^{2n}, \]
dengan setiap suku memuat salah satu dari $x,y$ berpangkat $ \geq n$, sehingga
setiap suku
termasuk $I$. Jadi
$(x+y)^{2n} \in I$, dan menurut definisi
terlihat bahwa $x+y \in \rad(I)$.
\end{proof} 

Pemetaan $I \to V(I)$ sesungguhnya hanya bergantung pada radikal
$I$. Memang, jika $I,J$ memiliki radikal yang sama, $\rad(I) =
\rad(J)$, maka $V(I) = V(J)$.
Memang, $V(I) = V(\rad(I)) = V(\rad(J)) = V(J)$ menurut:
\begin{lemma} 
Untuk setiap $I$, berlaku $V(I) = V(\rad(I))$.
\end{lemma} 
\begin{proof} 
Memang, $I \subset \rad(I)$ dan karena itu jelas bahwa $V(\rad(I))
\subset V(I)$. Kita harus menunjukkan
inklusi sebaliknya. Tepatnya, kita harus membuktikan:
\begin{quote}
Jika $\mathfrak{p} \supset I$, maka $\mathfrak{p} \supset
\rad(I).$
\end{quote}
Jadi andaikan $\mathfrak{p} \supset I$ prima dan $x \in
\rad(I)$; maka $x^n \in I \subset \mathfrak{p}$ untuk suatu $n$.
Namun, $\mathfrak{p}$ prima, sehingga jika suatu hasil kali elemen-elemen
termasuk
$\mathfrak{p}$, salah satu faktornya pun demikian. Karena $x^n = x \cdot x
\cdots x$, kita harus
memiliki $x \in \mathfrak{p}$. Jadi
\[ \rad(I) \subset \mathfrak{p},  \]
yang membuktikan pernyataan dalam kutipan dan dengan demikian lema tersebut.
\end{proof} 

Pernyataan ini memiliki suatu kebalikan:
\begin{proposition} 
Jika $V(I) = V(J)$, maka $\rad(I) = \rad(J)$. 
\end{proposition} 
Jadi dua ideal mendefinisikan subhimpunan tertutup yang sama jika dan hanya jika keduanya memiliki
radikal yang sama.
\begin{proof} 
Kita tuliskan suatu rumus bagi $\rad(I)$ yang segera menyiratkan
pernyataan tersebut.
\begin{lemma} \label{radprimescontaining} Untuk gelanggang komutatif $R$ dan suatu
ideal $I \subset
R$,
\[ \rad(I) = \bigcap_{\mathfrak{p} \supset I} \mathfrak{p}.  \]
\end{lemma} 
Dari sini, $V(I)$ menentukan $\rad(I)$. Karena itu, rumus ini
menyiratkan
proposisi tersebut.  
Sekarang kita buktikan lema itu:
\begin{proof} 
\begin{enumerate}
\item Kita tunjukkan $\rad(I) \subset \bigcap_{\mathfrak{p} \in V(I)}
\mathfrak{p} $. Khususnya,
hal ini berlaku jika kita menunjukkan bahwa apabila suatu ideal prima memuat
$I$, ia juga memuat $\rad(I)$; tetapi hal ini sudah kita
bahas di atas.  
\item Jika $x \notin \rad(I)$, kita akan menunjukkan bahwa terdapat suatu ideal
prima $\mathfrak{p}
\supset I$ yang tidak memuat $x$. Hal ini menyiratkan inklusi
sebaliknya dan
lema tersebut.  
\end{enumerate}


Kita ingin mencari $\mathfrak{p}$ yang tidak memuat $x$, atau secara lebih
umum tidak
memuat pangkat mana pun dari $x$. Khususnya, kita menginginkan
$\mathfrak{p} \cap \left\{1,
x, x^2 \dots, \right\} = \emptyset$. Himpunan $S = \left\{1, x,
\dots\right\}$ ini
tertutup secara multiplikatif, yakni memuat 1 dan tertutup
terhadap
hasil kali hingga. Saat ini, himpunan itu tidak beririsan dengan $I$; kita ingin mencari
suatu ideal
\emph{prima} yang memuat $I$ tetapi tetap tidak beririsan dengan $\left\{x^n, n
\geq 0\right\}$.


Secara lebih umum, kita akan membuktikan:

\begin{sublemma}\label{sublemmamultclosed}
Misalkan $S$ suatu himpunan tertutup secara multiplikatif dalam sembarang gelanggang $R$, dan misalkan
$I$ suatu ideal dengan $I \cap S =
\emptyset$. Terdapat suatu ideal prima $\mathfrak{p} \supset I$ yang
tidak
beririsan dengan $S$ (sesungguhnya, setiap ideal yang maksimal sehubungan dengan syarat
tidak beririsan dengan $S$ dapat digunakan).  
\end{sublemma}
Dengan kata lain, setiap ideal yang tidak bertemu $S$ dapat diperbesar menjadi suatu ideal
prima yang juga tidak bertemu $S$.
Ini sesungguhnya versi yang lebih canggih dari suatu argumen sebelumnya. Sebelumnya kita
telah menunjukkan bahwa setiap ideal yang tidak
memuat subhimpunan tertutup secara multiplikatif $\left\{1\right\}$
dapat
dimuat dalam suatu ideal prima yang tidak memuat $1$, dalam
\rref{anycontainedinmaximal}.

Perhatikan bahwa sublema tersebut jelas menyiratkan lema ketika diterapkan pada
$S =
\left\{1, x, \dots\right\}.$

\begin{proof}[Bukti sublema]
Misalkan $P = \left\{J: J \supset I, J \cap S = \emptyset \right\}$.
Maka $P$ adalah suatu
poset terhadap inklusi. Perhatikan bahwa $P \neq \emptyset$
karena $I \in P$. Selain itu,
untuk setiap subhimpunan tak kosong yang terurut linear dari $P$, gabungannya berada dalam
$P$ (yakni, terdapat suatu
batas atas).  
Kita dapat menggunakan Lema Zorn untuk memperoleh suatu elemen maksimal dari $P$. Elemen ini
adalah suatu
ideal $\mathfrak{p} \supset I$ dengan $\mathfrak{p} \cap S =
\emptyset$. Kita menyatakan
bahwa $\mathfrak{p}$ prima.

Pertama-tama, $1 \notin \mathfrak{p}$ karena $1 \in S$. Kita hanya perlu
memeriksa
bahwa jika $xy \in \mathfrak{p}$, maka $x \in \mathfrak{p}$ atau $y
\in
\mathfrak{p}$. Andaikan tidak demikian, sehingga $x,y \notin \mathfrak{p}$.
Maka $(x,\mathfrak{p}) \notin P$, atau
$\mathfrak{p}$ tidak maksimal. Hal yang sama berlaku bagi $(y,
\mathfrak{p})$.

Khususnya, kedua ideal yang lebih besar ini beririsan dengan
$S$. Artinya,
terdapat 
\[ a \in \mathfrak{p} , r \in R \quad \text{sedemikian sehingga}\quad a+rx
\in S \]
dan 
\[ b \in \mathfrak{p} , r' \in R \quad \text{sedemikian sehingga}\quad
b+r'y \in S .\]
Sekarang $S$ tertutup secara multiplikatif, maka kalikan $(a+rx)(b+r'y)
\in S$.
Kita memperoleh:
\[ ab + ar'y+brx+rr'xy \in S.  \]
Karena $a,b \in \mathfrak{p}$ dan $xy \in \mathfrak{p}$, semua
suku di atas berada dalam $\mathfrak{p}$, demikian pula jumlahnya. Namun, ini bertentangan dengan
$\mathfrak{p}
\cap S = \emptyset$. 
\end{proof}
\end{proof} 
\end{proof} 
Kesimpulan dari lema-lema sebelumnya adalah:
\begin{proposition}
Terdapat suatu bijeksi antara subhimpunan tertutup dari $\spec R$ dan
ideal-ideal radikal
$I \subset R$.
\end{proposition}

\subsection{Suatu pengamatan meta tentang ideal prima}

Dalam subbagian sebelumnya (\cref{sublemmamultclosed}), kita melihat
bahwa suatu ideal yang maksimal sehubungan dengan sifat tidak beririsan dengan suatu
subhimpunan tertutup secara multiplikatif adalah prima.
Ternyata hal serupa berlaku bagi banyak sifat ideal semacam ini.
Suatu metode umum untuk melihatnya dikembangkan dalam \cite{LaRe08}.
Dalam subbagian (opsional) ini, kita menyimpang sejenak untuk menjelaskan fenomena tersebut.

Jika $I$ suatu ideal dan $a \in R$, kita definisikan notasi
\[ (I:a) = \left\{ x\in R: xa \in I\right\} . \]
Secara lebih umum, jika $J$ suatu ideal, kita definisikan
\[ (I:J) = \left\{x \in R: xJ \subset I\right\} . \]

Misalkan $R$ suatu gelanggang dan $\mathcal{F}$ suatu koleksi ideal dari $R$. 
Kita tertarik pada syarat-syarat yang menjamin bahwa elemen-elemen maksimal
dari $\mathcal{F}$ bersifat \emph{prima}.
Sesungguhnya, kita akan melakukan yang sebaliknya: syarat berikut akan menjamin bahwa
ideal-ideal yang maksimal dalam hal \emph{tidak} termasuk $\mathcal{F}$ bersifat prima. 

\begin{definition} \label{okafamily}
Keluarga $\mathcal{F}$ disebut suatu \textbf{keluarga Oka} jika $R \in
\mathcal{F}$ (dengan $R$ dipandang sebagai ideal), dan jika $I \subset R$ suatu
ideal serta $(I:a), (I,a) \in \mathcal{F}$ (untuk suatu $a \in R$), maka $I \in
\mathcal{F}$.
\end{definition} 

\begin{example} \label{exm:okacard}
Mari kita mulai dengan suatu pengamatan sederhana. Jika $(I:a)$ dibangkitkan oleh
$a_1, \dots, a_n$ dan $(I,a)$ dibangkitkan oleh $a, b_1, \dots, b_m$ (dengan
tanpa mengurangi keumuman, kita dapat mengambil
$b_1, \dots, b_m \in I$), maka $I$
dibangkitkan oleh $aa_1, \dots, aa_n, b_1, \dots, b_m$.
Untuk melihatnya, perhatikan bahwa jika $x \in I$, maka $x \in (I,a)$ merupakan kombinasi
linear dari $\left\{a, b_1, \dots, b_m\right\}$, tetapi koefisien
$a$ harus
terletak dalam $(I:a)$.

Akibatnya, kita dapat menyimpulkan bahwa
keluarga ideal terbangkitkan hingga merupakan suatu keluarga Oka.
\end{example} 

\begin{example} 
Sekarang kita tunjukkan bahwa keluarga ideal \emph{prinsipal} merupakan suatu keluarga Oka.
Memang, andaikan $I \subset R$ suatu ideal, dan $(I,a)$ serta $(I:a)$ keduanya
prinsipal.
Kita dapat dengan mudah memeriksa bahwa
$(I:a) = (I: (I, a))$.
Dengan menetapkan $J = (I,a)$, kita memperoleh bahwa $J$ prinsipal dan $(I:J)$ juga demikian.
Namun, untuk \emph{setiap} ideal prinsipal $J$, dan untuk setiap ideal $I \subset J$,
\[ I = J (I: J)  \]
seperti yang mudah diperiksa. Jadi, dalam situasi kita, karena $J=(I,a)$ dan
$(I:J)$
prinsipal, maka $I$ prinsipal.
\end{example} 

\begin{proposition}[\cite{LaRe08}]\label{okathm} Jika $\mathcal{F}$ suatu keluarga Oka
dari
ideal-ideal, maka setiap elemen maksimal dari komplemen $\mathcal{F}$ bersifat prima.
\end{proposition} 
\begin{proof} 
Andaikan $I \notin \mathcal{F}$ maksimal sehubungan dengan 
sifat tidak termasuk $\mathcal{F}$,
tetapi $I$ tidak prima. Perhatikan bahwa $I \neq R$ menurut hipotesis.
Maka terdapat $a \in R$ sedemikian sehingga $(I:a), (I,a)$ keduanya memuat $I$ secara ketat,
sehingga keduanya harus termasuk $\mathcal{F}$.
Memang, kita dapat menemukan $a,b \in R - I$ dengan $ab \in I$; akibatnya, $(I,a)
\neq I$, dan $(I:a)$ memuat $b \notin I$.

Menurut syarat Oka, kita memiliki $I \in
\mathcal{F}$, suatu kontradiksi.
\end{proof} 

\begin{corollary}[Cohen] \label{primenoetherian}
Jika setiap ideal prima $R$ terbangkitkan hingga, maka setiap ideal $R$
terbangkitkan hingga.\footnote{Kelak kita akan mengatakan bahwa $R$ bersifat \emph{noetherian}.} 
\end{corollary} 

\begin{proof} 
Andaikan terdapat ideal-ideal $I \subset R$ yang tidak terbangkitkan
hingga.
Gabungan suatu rantai terurut total $\left\{I_\alpha\right\}$ dari ideal-ideal yang
tidak terbangkitkan hingga juga tidak
terbangkitkan hingga; memang, jika $I = \bigcup I_\alpha$ dibangkitkan oleh $a_1, \dots,
a_n$, maka semua pembangkit tersebut akan termasuk suatu $I_\alpha $ dan akibatnya
membangkitkannya.

Menurut Lema Zorn, terdapat suatu ideal yang maksimal sehubungan dengan sifat tidak
terbangkitkan hingga. Namun, menurut \rref{okathm}, ideal ini harus
prima (karena keluarga ideal terbangkitkan hingga merupakan suatu keluarga Oka). Hal ini bertentangan dengan hipotesis.
\end{proof} 

\begin{corollary} 
Jika setiap ideal prima $R$ prinsipal, maka setiap ideal $R$ prinsipal.
\end{corollary} 
\begin{proof} 
Pernyataan ini dibuktikan dengan cara yang sama.
\end{proof} 

\begin{exercise} 
Andaikan setiap ideal prima tak nol dalam $R$ memuat suatu elemen yang bukan pembagi nol. Maka $R$
merupakan suatu domain. (Petunjuk: tinjau himpunan $S$ dari elemen-elemen yang bukan pembagi nol, dan tunjukkan bahwa
setiap ideal yang maksimal sehubungan dengan sifat tidak beririsan dengan $S$ bersifat prima. Dengan demikian, $(0)$
prima.)
\end{exercise} 


\begin{remark}
\label{remark-cohen-bound-cardinality}
Misalkan $R$ suatu gelanggang. Misalkan $\kappa$ suatu kardinal tak hingga.
Dengan menerapkan
\rref{exm:okacard} dan
\rref{okathm},
kita melihat bahwa setiap ideal yang maksimal sehubungan dengan sifat tidak
terbangkitkan oleh $\kappa$ elemen bersifat prima. Hasil ini tidak begitu
berguna karena terdapat suatu gelanggang yang setiap ideal prima
dari $R$ dapat dibangkitkan oleh $\aleph_0$ elemen, tetapi suatu
idealnya tidak dapat. Tepatnya, misalkan $k$ suatu medan, $T$ suatu himpunan yang
kardinalitasnya lebih besar daripada $\aleph_0$, dan tetapkan
\[ R = k[\{x_n\}_{n \geq 1}, \{z_{t, n}\}_{t \in T, n \geq 0}]/
(x_n^2, z_{t, n}^2, x_n z_{t, n} - z_{t, n - 1}) \]
Ini adalah gelanggang lokal dengan ideal prima tunggal
$\mathfrak m = (x_n)$. Namun, ideal $(z_{t, n})$ tidak dapat
dibangkitkan oleh terhitung banyak elemen.
\end{remark}

\subsection{Sifat funktorial $\spec$}
 Konstruksi $R \to \spec R$ bersifat funktorial dalam $R$ dalam arti
 kontravarian. Artinya, jika $f: R \to R'$, terdapat suatu
pemetaan kontinu $\spec
R' \to \spec R$. Pemetaan ini mengirim $\mathfrak{p} \subset R'$ ke
$f^{-1}(\mathfrak{p}) \subset R$, yang mudah dilihat sebagai suatu ideal
prima
dalam $R$. Sebut pemetaan ini $F: \spec R' \to \spec R$. Sejauh ini, kita telah
melihat bahwa
$\spec R$ menginduksi suatu funktor kontravarian dari $\mathbf{Rings}
\to \mathbf{Sets}$.

\begin{exercise} 
Suatu funktor kontravarian $F: \mathcal{C} \to \mathbf{Sets}$ (untuk
suatu kategori
$\mathcal{C}$) disebut \textbf{terrepresentasikan} jika
isomorfik secara alami
dengan suatu funktor berbentuk $X \to \hom(X, X_0)$ untuk suatu $X_0 \in
\mathcal{C}$,
atau, secara ekuivalen, jika funktor kovarian terinduksi pada
$\mathcal{C}^{\mathrm{op}}$ terkorepresentasikan. 

Funktor $R \to \spec R $ tidak terrepresentasikan. (Petunjuk:
memang, suatu funktor terrepresentasikan
harus mengirim objek awal ke suatu himpunan satu titik.)
\end{exercise} 

Selanjutnya, kita periksa bahwa morfisme pada $\spec$ yang diinduksi oleh suatu
homomorfisme gelanggang memang merupakan pemetaan \emph{kontinu} antarruang
topologis.

\begin{proposition} 
$\spec $ menginduksi suatu funktor kontravarian dari $\mathbf{Rings}$
ke kategori
$\mathbf{Top}$ ruang topologis.
\end{proposition} 
\begin{proof} Misalkan $f : R \to R'$. 
Kita perlu memeriksa bahwa pemetaan $ \spec R' \to \spec R$ ini, yang
kita sebut $F$,
kontinu.
Artinya, kita harus memeriksa bahwa $F^{-1}$ mengirim subhimpunan tertutup
dari $\spec R$ ke subhimpunan tertutup dari $\spec R'$.  

Lebih tepatnya, jika $I \subset
R$ dan kita mengambil prapeta $F^{-1}(V(I)) \subset \spec
R'$, maka himpunan itu tepat
himpunan tertutup $V(f(I))$. Hal ini sebaiknya diserahkan kepada pembaca, tetapi
berikut alasannya. Jika $\mathfrak{p} \in \spec R'$, maka
$F(\mathfrak{p}) = f^{-1}(\mathfrak{p})
\supset I$ jika dan hanya jika $\mathfrak{p} \supset f(I)$. Jadi
$F(\mathfrak{p}) \in
V(I)$ jika dan hanya jika $\mathfrak{p} \in V(f(I))$.

\end{proof} 



\begin{example} 
Misalkan $R$ suatu gelanggang komutatif, $I \subset R$ suatu ideal, dan $f: R \to
R/I$. Terdapat suatu pemetaan
ruang topologis
\[ F: \spec (R/I) \to \spec R  .\]
Pemetaan ini adalah suatu embedding tertutup dengan citra $V(I)$. Sebagian besar
pernyataan ini berlaku karena
terdapat suatu bijeksi antara ideal-ideal $R$ yang memuat $I$ dan
ideal-ideal $R/I$, dan bijeksi ini mempertahankan keprimaan.

\begin{exercise} 
Tunjukkan bahwa pemetaan $\spec R/I \to \spec R$ ini memang suatu
homeomorfisme dari $\spec R/I
\to V(I)$.  
\end{exercise} 
\end{example} 


\subsection{Suatu basis bagi topologi Zariski}
Dalam bagian sebelumnya, kita membahas topologi Zariski.
Jika $R$ suatu
gelanggang komutatif, ingat bahwa $\spec R$ didefinisikan sebagai
koleksi
ideal prima dalam $R$. Himpunan ini memiliki topologi yang himpunan tertutupnya
adalah himpunan-himpunan
berbentuk
\[ V(I) = \left\{\mathfrak{p} \in \spec R: \mathfrak{p} \supset
I\right\} . \]
Ada cara lain untuk mendeskripsikan topologi Zariski dalam kaitannya dengan
himpunan
\emph{terbuka}.  

\begin{definition} 
Jika $f \in R$, tetapkan 
\[ U_f = \left\{\mathfrak{p}: f \notin \mathfrak{p}\right\}  \]
sehingga $U_f$ adalah subhimpunan $\spec R$ yang terdiri dari ideal-ideal prima
yang tidak memuat
$f$. Himpunan ini adalah komplemen $V((f))$, sehingga terbuka.  
\end{definition} 

\begin{proposition} 
Himpunan-himpunan $U_f$ membentuk suatu basis bagi topologi Zariski. 
\end{proposition} 

\begin{proof} 
Andaikan $U \subset \spec R$ terbuka. Kita menyatakan bahwa $U$ merupakan suatu
gabungan himpunan-himpunan
terbuka basis $U_f$. 

Sekarang $U = \spec R - V(I)$ untuk suatu ideal $I$. Maka
\[ U = \bigcup_{f \in I} U_f  \]
karena jika suatu ideal tidak berada dalam $V(I)$, maka ia gagal memuat
suatu $f \in I$,
yakni berada dalam $U_f$ untuk $f$ tersebut. Sebagai alternatif, kita dapat mengambil
komplemen, sehingga
pernyataan di atas menjadi
\[ V(I) = \bigcap_{f \in I} V((f))  \]
yang jelas.
\end{proof} 

Himpunan-himpunan terbuka basis memiliki sifat-sifat yang baik.
\begin{enumerate}
\item $U_1 = \spec R$ karena ideal prima tidak diperbolehkan memuat
elemen
unit. 
\item $U_0 = \emptyset$ karena setiap ideal prima memuat $0$.
\item $U_{fg} = U_f \cap U_g$ karena $fg$ terletak dalam suatu ideal prima
$\mathfrak{p}$ jika dan hanya jika salah satu dari $f,g$ demikian.
\end{enumerate}

Sekarang mari kita deskripsikan hubungan topologi Zariski dengan
lokalisasi.
Misalkan $R$ suatu gelanggang dan $f \in R$. Tinjau $S = \left\{1, f, f^2,
\dots
\right\}$; ini suatu subhimpunan tertutup secara multiplikatif. Sebelumnya,
kita mendefinisikan
$S^{-1}R$.

\begin{definition} 
Jika $S$ adalah himpunan pangkat $f$, kita menulis $R_f$ atau $R[f^{-1}]$ untuk
lokalisasi $S^{-1}R$. 
\end{definition} 

Terdapat suatu pemetaan $\phi: R \to R[f^{-1}]$ dan pemetaan yang bersesuaian
\[ \spec R[f^{-1}] \to \spec R  \]
yang mengirim suatu ideal prima $\mathfrak{p} \subset R[f^{-1}]$ ke
$\phi^{-1}(\mathfrak{p})$.

\begin{proposition} 
Pemetaan ini menginduksi suatu homeomorfisme dari $\spec R[f^{-1}]$ ke $U_f
\subset \spec
R$. 
\end{proposition} 
Jadi, jika suatu elemen dalam gelanggang komutatif diinverskan, yang diperoleh hanyalah suatu subhimpunan
terbuka dari $\spec$. Inilah alasan operasi tersebut disebut lokalisasi:
kita
membatasi perhatian pada
suatu subhimpunan terbuka pada taraf $\spec $ ketika menginverskan sesuatu.

\begin{proof}
Pembaca dianjurkan mengerjakan sendiri pembuktian ini.

\begin{enumerate}
\item 
Pertama, kita tunjukkan bahwa $\spec R[f^{-1}] \to \spec R$ mendarat dalam
$U_f$. Jika
$\mathfrak{p} \subset R[f^{-1}]$, maka kita harus menunjukkan bahwa
prapeta
$\phi^{-1}(\mathfrak{p})$ tidak dapat memuat $f$. Jika sebaliknya, hal itu
akan menyiratkan bahwa
$\phi(f) \in \mathfrak{p}$; namun, $\phi(f)$ invertibel,
dan kemudian
$\mathfrak{p}$ akan sama dengan $(1)$.  
\item Mari kita tunjukkan bahwa pemetaan tersebut surjektif ke $U_f$. Jika
$\mathfrak{p} \subset R$ suatu ideal
prima yang tidak memuat $f$, yakni $\mathfrak{p} \in U_f$, kita ingin
mengkonstruksi suatu
ideal prima yang bersesuaian dalam gelanggang $R[f^{-1}]$ dengan prapeta
$\mathfrak{p}$.

Misalkan $\mathfrak{p}[f^{-1}]$ koleksi semua pecahan
\[ \{\frac{x}{f^n}, x \in \mathfrak{p}\} \subset R[f^{-1}],  \]
yang jelas merupakan suatu ideal. Perhatikan bahwa apakah pembilang berada
dalam
$\mathfrak{p}$ \textbf{tidak bergantung} pada
pecahan representatif $\frac{x}{f^n}$ yang digunakan.\footnote{Andaikan
$\frac{x}{f^n} =
\frac{y}{f^k}$ untuk $y \in \mathfrak{p}$. Maka terdapat $N$ sedemikian
sehingga
$f^N(f^k x - f^n y) = 0 \in \mathfrak{p}$; karena $y \in
\mathfrak{p}$ dan $f
\notin \mathfrak{p}$, diperoleh $x \in \mathfrak{p}$.}
Sesungguhnya, $\mathfrak{p}[f^{-1}]$ suatu ideal prima. Memang,
andaikan
\[  \frac{a}{f^m} \frac{b}{f^n} \in \mathfrak{p}[f^{-1}] .\]
Maka $\frac{ab}{f^{m+n}}$ termasuk ideal ini, yang berarti $ab
\in
\mathfrak{p}$; jadi salah satu dari $a,b \in \mathfrak{p}$, dan salah satu
dari kedua pecahan
$\frac{a}{f^m}, \frac{b}{f^n}$ termasuk
$\mathfrak{p}[f^{-1}]$. Selain itu, $1/1
\notin \mathfrak{p}[f^{-1}]$.

Jelas bahwa prapeta $\mathfrak{p}[f^{-1}]$ adalah
$\mathfrak{p}$,
karena citra $x \in R$ adalah $x/1$, dan elemen ini termasuk
$\mathfrak{p}[f^{-1}]$ tepat ketika $x \in \mathfrak{p}$.
\item Pemetaan $\spec R[f^{-1}] \to \spec R$ injektif.
Andaikan
$\mathfrak{p}, \mathfrak{p'}$ merupakan ideal prima dalam
lokalisasi tersebut dan
prapetanya sama.  
Kita harus menunjukkan bahwa $\mathfrak{p} = \mathfrak{p'}$.

Andaikan $\frac{x}{f^n} \in \mathfrak{p}$. Maka $x/1 \in
\mathfrak{p}$, sehingga $x
\in \phi^{-1}(\mathfrak{p}) = \phi^{-1}(\mathfrak{p}')$. Ini
berarti bahwa $x/1
\in \mathfrak{p}'$, sehingga 
$\frac{x}{f^n} \in \mathfrak{p}'$ juga. Jadi setiap pecahan yang
termasuk
$\mathfrak{p}$ juga termasuk $\mathfrak{p}'$. Berdasarkan simetri, kedua ideal itu
harus
sama.  
\item Sekarang kita mengetahui bahwa pemetaan $\psi: \spec R[f^{-1}] \to U_f$
merupakan bijeksi
kontinu. Tinggal diperiksa bahwa pemetaan itu suatu homeomorfisme. Kita akan
menunjukkan bahwa pemetaan itu
terbuka.  
Khususnya, kita harus menunjukkan bahwa suatu himpunan terbuka basis di ruas kiri
dipetakan
ke himpunan terbuka di ruas kanan.
Jika $y/f^n \in R[f^{-1}]$, kita harus menunjukkan bahwa $U_{y/f^n}
\subset \spec
R[f^{-1}]$ memiliki citra terbuka di bawah $\psi$. Bahkan, kita akan menunjukkan himpunan
terbuka mana yang menjadi citranya.

Kita menyatakan bahwa
\[ \psi(U_{y/f^n}) = U_{fy} \subset \spec R.  \]
Untuk melihatnya, $\mathfrak{p}$ termasuk dalam $U_{f/y^n}$. Ini
berarti bahwa
$\mathfrak{p}$ tidak memuat $y/f^n$. Khususnya,
$\mathfrak{p}$ tidak
memuat kelipatan $yf/1$. Jadi $\psi(\mathfrak{p})$ tidak
memuat $yf$.
Hal ini membuktikan inklusi $\subset$.  

\item 
Untuk menyelesaikan pembuktian pernyataan tersebut, dan
hasilnya, kita harus menunjukkan bahwa jika $\mathfrak{p} \subset \spec
R[f^{-1}]$ dan
$\psi(\mathfrak{p}) = \phi^{-1}(\mathfrak{p}) \in U_{fy}$, maka
$y/f^n$ tidak
termasuk $\mathfrak{p}$. (Ini sah-sah saja karena
kita memiliki suatu bijeksi.) Namun, hipotesis mengakibatkan $fy \notin
\phi^{-1}(\mathfrak{p})$, sehingga $fy/1 \notin \mathfrak{p}$.
Membagi dengan $f^{n+1}$
mengakibatkan bahwa
\[ y/f^{n} \notin \mathfrak{p}  \]
dan $\mathfrak{p} \in U_{f/y^n}$. 
\end{enumerate}
\end{proof}

Jika $\spec R$ suatu ruang, dan $f$ dipandang sebagai suatu ``fungsi'' yang
didefinisikan pada
$\spec R$, maka ruang $U_f$ harus dipandang sebagai himpunan
titik tempat $f$
``tidak lenyap'' atau ``invertibel''.
Memandang gelanggang melalui spektrumnya merupakan gagasan yang sangat berguna.
Kita akan kembali menggunakannya jika sesuai.  

\begin{remark} 
Konstruksi $R \to R[f^{-1}]$ yang dibahas di atas merupakan suatu
contoh
lokalisasi. Secara lebih umum, kita mendefinisikan $S^{-1}R$ untuk $S
\subset R$
yang tertutup secara multiplikatif. Dengan demikian, kita dapat mendefinisikan pemetaan
\( \spec S^{-1}R \to \spec R . \)
Untuk memahami $S^{-1}R$, mungkin membantu untuk mencatat bahwa
\[ \varinjlim_{f \in S} R[f^{-1}]  \]
merupakan limit langsung gelanggang-gelanggang yang menginverskan semakin banyak elemen.

Sebagai contoh, tinjau $S = R - \mathfrak{p}$ untuk suatu ideal prima
$\mathfrak{p}$, dan demi
kesederhanaan andaikan $R$ terhitung. Kita dapat menulis $S =
S_0 \cup S_1 \cup \dots$, dengan setiap $S_k$ dibangkitkan oleh
sejumlah hingga
elemen $f_0, \dots, f_k$. Maka $R_{\mathfrak{p}} = \varinjlim
S_k^{-1} R$.
Jadi kita memiliki
\[ S^{-1}R = \varinjlim_k R[f_0^{-1} , f_1^{-1}, \dots, f_k^{-1}
] = \varinjlim
R[(f_0\dots f_k)^{-1}]. \]
Elemen-elemen yang kita inverskan dalam konstruksi ini tepat merupakan elemen-elemen yang
tidak
termasuk $\mathfrak{p}$, atau tempat ``fungsi-fungsi tidak lenyap''.

Gagasan geometrisnya adalah
bahwa untuk mengkonstruksi $\spec S^{-1}R = \spec R_{\mathfrak{p}}$, kita
terus membuang
dari $\spec R$ lokus-lokus lenyap berbagai fungsi yang
tidak
beririsan dengan $\mathfrak{p}$. Pada akhirnya, kita tidak membatasi diri pada suatu
himpunan terbuka, melainkan
pada irisan himpunan-himpunan tersebut.
\end{remark} 
\begin{exercise} \label{semilocal} 
Katakan bahwa $R$ \emph{semilokal} jika memiliki berhingga banyak ideal maksimal.
Misalkan $\mathfrak{p}_1$, \dots, $\mathfrak{p}_n\subset R$ ideal-ideal prima. Komplemen
gabungannya, $S=R\smallsetminus \bigcup \mathfrak{p}_i$, tertutup terhadap
perkalian, sehingga kita dapat
melokalkan. $R[S^{-1}] = R_S$ disebut \emph{semilokalisasi}
 \index{semi-localization} dari $R$ di $\mathfrak{p}_i$.
 
Hasil semilokalisasi selalu semilokal. Untuk melihatnya, ingat bahwa
 ideal-ideal
 dalam $R_S$ berkorespondensi bijektif dengan ideal-ideal dalam $R$ yang termuat dalam $\bigcup
 \mathfrak{p}_i$. Sekarang gunakan penghindaran prima.
\end{exercise}

\begin{definition}
Untuk suatu modul-$R$ terbangkitkan hingga $M$, definisikan $\mu_R(M)$ sebagai jumlah terkecil
elemen
   yang dapat membangkitkan $M$.
 \end{definition}
Ini tidak sama dengan kardinalitas suatu himpunan pembangkit minimal. Sebagai
contoh, 2
dan 3 merupakan suatu himpunan pembangkit minimal bagi $\mathbb{Z}$ atas dirinya sendiri, tetapi
$\mu_\mathbb{Z} (\mathbb{Z}) =1$.

 \begin{theorem}
Misalkan $R$ semilokal dengan ideal-ideal maksimal $\mathfrak{m}_1,\dots,
\mathfrak{m}_n$. Misalkan $k_i = R/\mathfrak{m}_i$. Maka
   \[
     \mu_R(M) = \max_i \{\dim_{k_i} M/\mathfrak{m}_i M\}
   \]
 \end{theorem}
\begin{proof} 
\begin{quote}\small\textbf{Jembatan edisi (asli).}
Ketaksamaan ``$\geq$'' diperoleh dengan mereduksi setiap himpunan pembangkit modulo setiap
$\mathfrak m_i$. Sebaliknya, tetapkan $d=\max_i\dim_{k_i}(M/\mathfrak m_iM)$
dan lengkapi suatu basis dari setiap $M/\mathfrak m_iM$ menjadi tupel-$d$. Teorema
Sisa Cina membuat pemetaan
\[
  M\longrightarrow\bigoplus_i M/\mathfrak m_iM
\]
surjektif, jadi pilih $x_1,\ldots,x_d\in M$ yang citranya adalah tupel-tupel tersebut.
Jika $N=\sum_jRx_j$, maka $(M/N)_{\mathfrak m_i}=0$ untuk setiap ideal maksimal dari
gelanggang semilokal $R$, menurut Lema Nakayama. Jadi $M/N=0$, dan elemen-elemen $x_j$
membangkitkan $M$.\end{quote}
\end{proof} 
 
